\chapter{Review of analytic methods for solving ODEs and PDEs}

Here we recall analytic tricks used to solve ODEs and PDEs in closed form.
These tricks only work for narrow, special classes of problems --
which is precisely why the numerical methods developed throughout the
rest of this course are needed -- but they remain indispensable:
closed-form solutions are our main source of test problems for
checking numerical algorithms, and several of the ideas below
(eigen-decompositions, separation of variables, Fourier and Green's
function ideas) reappear directly inside the numerical methods
themselves. As before, this is a brief review; see the classic texts
\cite{boyce2017elementary,strauss2007partial} for a much more detailed
treatment.

\section{Solving ODEs}

Throughout this section we consider ODEs of the general form
\eqref{eq:first_order_ode} or \eqref{eq:high_order_ode} from
\Cref{ch:Lec1-intro}, together with the IVP and BVP framework of
\Cref{def:ivp,def:bvp}.

\subsection{General and particular solutions}

An ODE of order $n$ typically has a \emph{general solution}: a family
of solutions depending on $n$ arbitrary constants, one for each order
of differentiation. Fixing the constants (e.g., by imposing initial or
boundary conditions as in \Cref{def:ivp,def:bvp}) singles out one
member of the family, called a \emph{particular solution}. We already
saw an instance of this in \Cref{ex:harmonic_oscillator}: the general
solution $u(t) = A\cos(\omega t) + B \sin(\omega t)$ of the harmonic
oscillator becomes the particular solution
\eqref{eq:harmonic_oscillator_solution} once we impose the initial
condition $u(0) = u_0$, $u'(0) = v_0$.

For \emph{linear} ODEs there is a very useful way to organize this
search, based on the following superposition principle.

\begin{proposition}[Structure of solutions to linear ODEs]
    Consider a linear ODE $Lu = g$, where $L$ is a linear differential
    operator (e.g., $Lu = u' + p(t) u$). If $u_p$ is any one particular
    solution of $Lu=g$ and $u_h$ is the general solution of the
    associated homogeneous equation $Lu=0$, then $u = u_h + u_p$ is the
    general solution of $Lu=g$.
\end{proposition}
This reduces solving an inhomogeneous linear ODE to two separate, often
easier, tasks: find the general homogeneous solution, and find just
one particular solution of the full equation.

\begin{example}
    Consider $u' - u = e^{2t}$. The homogeneous equation $u'-u=0$ has
    general solution $u_h(t) = Ce^t$. For a particular solution we
    guess $u_p(t) = ae^{2t}$; substituting gives $2ae^{2t} - ae^{2t} =
    e^{2t}$, so $a=1$. Hence the general solution is
    \begin{equation*}
        u(t) = Ce^t + e^{2t}.
    \end{equation*}
\end{example}

\subsection{Separation of variables}

Separation of variables applies to first-order ODEs that can be
written in the form
\begin{equation*}
    \frac{du}{dt} = f(t) g(u).
\end{equation*}
Formally separating the variables and integrating both sides,
\begin{equation*}
    \int \frac{du}{g(u)} = \int f(t) \, dt,
\end{equation*}
gives an implicit (and often explicit) relation between $u$ and $t$
involving one arbitrary constant, which is then fixed by an initial
condition.

\begin{example}[Logistic growth]
    Consider the logistic equation $du/dt = ku(1-u)$, which models,
    e.g., a population growing toward a carrying capacity of $1$.
    Separating variables and using partial fractions,
    \begin{equation*}
        \int \frac{du}{u(1-u)} = \int \left( \frac1u + \frac{1}{1-u}
        \right) du = \ln \left| \frac{u}{1-u} \right| = kt + C.
    \end{equation*}
    Solving for $u$ and imposing $u(0) = u_0$ gives
    \begin{equation*}
        u(t) = \frac{u_0 e^{kt}}{1 - u_0 + u_0 e^{kt}},
    \end{equation*}
    which indeed satisfies $u(t) \to 1$ as $t \to \infty$ for any
    $u_0 \in (0,1)$.
\end{example}

\subsection{Integrating factors}

The integrating factor method solves general linear first-order ODEs
\begin{equation*}
    \frac{du}{dt} + p(t) u = q(t).
\end{equation*}
Multiplying both sides by $\mu(t) := \exp\left( \int p(t) \, dt
\right)$ turns the left-hand side into a total derivative,
$\frac{d}{dt}(\mu u) = \mu q$, i.e., we have multiplied by exactly the
factor needed to make the left-hand side the derivative of a product.
Integrating both sides then gives
\begin{equation*}
    u(t) = \frac{1}{\mu(t)} \left( \int \mu(t) q(t) \, dt + C \right).
\end{equation*}

\begin{example}[Newton's law of cooling]
    An object at temperature $u(t)$ in an environment at constant
    temperature $T_{\mathrm{env}}$ cools (or warms) at a rate
    proportional to the temperature difference, $du/dt = -k(u -
    T_{\mathrm{env}})$, i.e., $u' + ku = kT_{\mathrm{env}}$. The
    integrating factor is $\mu(t) = e^{kt}$, and integrating
    $\frac{d}{dt}(e^{kt}u) = kT_{\mathrm{env}} e^{kt}$ gives $e^{kt}
    u(t) = T_{\mathrm{env}} e^{kt} + C$. Imposing $u(0) = u_0$ fixes
    $C = u_0 - T_{\mathrm{env}}$, so
    \begin{equation*}
        u(t) = T_{\mathrm{env}} + (u_0 - T_{\mathrm{env}}) e^{-kt},
    \end{equation*}
    which decays exponentially toward $T_{\mathrm{env}}$, as expected.
\end{example}

\subsection{Exact solutions}

An ODE written in the differential form
\begin{equation*}
    M(t,u) + N(t,u) \frac{du}{dt} = 0
\end{equation*}
is called \emph{exact} if $\partial M/\partial u = \partial N/\partial
t$ (recall \Cref{lem:mixed-partials}). In that case there exists a
potential function $\Psi(t,u)$ with $\partial \Psi/\partial t = M$ and
$\partial \Psi/\partial u = N$, and the solution is given implicitly by
the level sets $\Psi(t,u) = C$. The potential is found by integrating
$M$ in $t$ (treating $u$ as a constant) and fixing the resulting
integration ``constant'' (a function of $u$) by matching against $N$.

\begin{example}
    Consider $(2tu + 3) + (t^2-1) \dfrac{du}{dt} = 0$, so $M = 2tu+3$
    and $N = t^2-1$. Since $\partial M/\partial u = 2t = \partial
    N/\partial t$, the equation is exact. Integrating $M$ in $t$ gives
    $\Psi(t,u) = t^2 u + 3t + h(u)$ for some function $h$; matching
    $\partial \Psi/\partial u = t^2 + h'(u)$ with $N = t^2-1$ forces
    $h(u) = -u$. Hence the solution is given implicitly by
    \begin{equation*}
        t^2 u + 3t - u = C.
    \end{equation*}
\end{example}

\subsection{Substitutions}

Many nonlinear ODEs that are not directly separable, exact, or linear
can be reduced to one of the previous cases by a clever change of
variable. Two classic families are:
\begin{itemize}
    \item \emph{Homogeneous equations} $du/dt = F(u/t)$: the
    substitution $v = u/t$ (so $u = tv$, $u' = v + tv'$) turns this
    into the separable equation $t \, dv/dt = F(v) - v$.
    \item \emph{Bernoulli equations} $u' + p(t) u = q(t) u^n$ for $n
    \neq 0,1$: the substitution $v = u^{1-n}$ turns this into the
    linear equation $v' + (1-n)p(t) v = (1-n) q(t)$, solvable by an
    integrating factor.
\end{itemize}

\begin{example}
    Consider the Bernoulli equation $u' - u = -u^2$ (a rewriting of the
    logistic equation from above, with $n=2$). Substituting $v =
    u^{-1}$, so $v' = -u^{-2} u'$, and dividing the original equation
    by $-u^2$ gives $-u^{-2}u' + u^{-1} = 1$, i.e.,
    \begin{equation*}
        v' + v = 1,
    \end{equation*}
    a linear equation solvable by an integrating factor as in the
    previous section.
\end{example}

\subsection{Transformations and series}

The following techniques all work by transforming the original ODE (or
its solution) into a different representation in which the problem
becomes algebraic or otherwise easier to solve.

\paragraph{Power series and the Frobenius method.}
If $u$ is expected to be analytic near $t=0$, we can posit a power
series ansatz $u(t) = \sum_{n=0}^\infty a_n t^n$, substitute into the
ODE, and match powers of $t$ to obtain a recurrence relation for the
coefficients $a_n$.

\begin{example}
    For $u' = u$, writing $u = \sum_{n\ge0} a_n t^n$ gives $u' =
    \sum_{n\ge1} n a_n t^{n-1} = \sum_{n\ge0} (n+1) a_{n+1} t^n$.
    Matching this with $u = \sum_{n\ge0} a_n t^n$ term by term gives the
    recurrence $a_{n+1} = a_n/(n+1)$, so $a_n = a_0/n!$ and $u(t) = a_0
    \sum_{n\ge0} t^n/n! = a_0 e^t$, recovering the familiar solution.
\end{example}
When the ODE has a singular point (a coefficient vanishing at the
expansion point), the plain power series ansatz above may fail to
produce a solution; the Frobenius method generalizes it by allowing a
leading factor $t^r$ for a possibly non-integer exponent $r$, and is
the standard tool for equations such as Bessel's equation.

\paragraph{Laplace transform.}
The Laplace transform of $u$ is $U(s) := \mathcal{L}\{u\}(s) =
\int_0^\infty e^{-st} u(t) \, dt$. Its key property is that it turns
differentiation into multiplication: $\mathcal{L}\{u'\}(s) = sU(s) -
u(0)$. Applying it to a linear ODE with constant coefficients turns it
into an \emph{algebraic} equation for $U(s)$, which can be solved and
then inverted (typically via a table of transforms) to recover $u(t)$.

\begin{example}
    For $u' + u = 0$, $u(0)=1$, taking the Laplace transform gives
    $sU(s) - 1 + U(s) = 0$, so $U(s) = 1/(s+1)$, which we recognize as
    the transform of $u(t) = e^{-t}$.
\end{example}

\paragraph{Fourier transform.}
The Fourier transform $\hat u(\omega) := \int_{-\infty}^\infty u(x)
e^{-i\omega x} \, dx$ plays a similar role but is most useful for
problems posed on the whole real line (rather than $t \ge 0$), and
turns differentiation in $x$ into multiplication by $i\omega$. We will
use it below to solve the heat equation on the whole real line.

\paragraph{Fourier series.}
If instead $u$ is defined on a bounded interval, say $[0,L]$, we can
expand it in the orthogonal basis of sines and cosines (or complex
exponentials) with period related to $L$,
\begin{equation*}
    u(x) = \frac{a_0}{2} + \sum_{n=1}^\infty \left( a_n \cos\left(
    \frac{2\pi n x}{L} \right) + b_n \sin\left( \frac{2\pi n x}{L}
    \right) \right),
\end{equation*}
with the coefficients $a_n, b_n$ computed using the orthogonality of
sines and cosines on $[0,L]$. We will see below that Fourier (sine or
cosine) series arise naturally as eigenfunction expansions for PDEs
with boundary conditions on $[0,L]$.

\paragraph{Green's functions.}
For a linear differential operator $L$ (in $t$, or in $\xb$ for a
PDE), the Green's function $G$ solves $LG = \delta$, where $\delta$ is
the Dirac delta ``function'' (an idealized unit impulse), together with
the relevant boundary or initial conditions. By linearity, the solution
of $Lu = f$ for a general right-hand side $f$ is then obtained by
superposition,
\begin{equation*}
    u = \int G(\cdot, \tau) f(\tau) \, d\tau.
\end{equation*}

\begin{example}
    \label{ex:greens-function-poisson}
    Consider the 1D Poisson problem from \Cref{ex:poisson-bvp} with
    homogeneous boundary conditions, $-u'' = f$ on $(0,1)$, $u(0) =
    u(1) = 0$. Its Green's function is
    \begin{equation*}
        G(x,\xi) = \begin{cases}
            x(1-\xi), & 0 \le x \le \xi, \\
            \xi(1-x), & \xi \le x \le 1,
        \end{cases}
    \end{equation*}
    which one can check directly solves $-\partial^2 G/\partial x^2 =
    \delta(x-\xi)$ in $x$ (it is linear in $x$ on each side of $\xi$,
    continuous, and has a unit jump in slope at $x=\xi$) along with the
    boundary conditions. The solution of the general problem is then
    $u(x) = \int_0^1 G(x,\xi) f(\xi) \, d\xi$.
\end{example}

\subsection{Solving systems}

For linear systems with constant coefficients, $d\ub/dt = A\ub$ with
$A \in \R^{d \times d}$ (a special case of \Cref{def:ivp}), the natural
tool is the eigen-decomposition of $A$ reviewed in
\Cref{ch:Lec0-linalg}. If $A$ is diagonalizable, $A = Q\Lambda
Q^{-1}$, then the general solution is
\begin{equation*}
    \ub(t) = Qe^{\Lambda t} Q^{-1} \ub(0) = \sum_{j=1}^d c_j
    e^{\lambda_j t} \qb_j,
\end{equation*}
where $\qb_j$ are the columns of $Q$ (the eigenvectors of $A$),
$\lambda_j$ the corresponding eigenvalues, and $c_j$ the coefficients
of $\ub(0)$ in the eigenvector basis; this follows directly from the
matrix exponential identity reviewed in \Cref{ch:Lec-1-calculus}. When
$A$ is not diagonalizable, one instead works with its Jordan form,
which introduces extra polynomial-in-$t$ terms alongside the
exponentials; we will not need the details of this construction in
this course.

\begin{example}
    The harmonic oscillator \eqref{eq:mass_spring} can be written as
    the first-order system $\ub' = A\ub$ with $\ub = (u,u')$, following
    the standard trick of introducing the velocity as an auxiliary
    variable (cf.\ \Cref{def:ivp}), where
    \begin{equation*}
        A = \begin{bmatrix} 0 & 1 \\ -\omega^2 & 0 \end{bmatrix}.
    \end{equation*}
    The eigenvalues of $A$ are $\lambda = \pm i\omega$ with eigenvectors
    $\qb_\pm = (1, \pm i \omega)$, so the eigen-decomposition gives
    solutions of the form $u(t) = c_+ e^{i\omega t} + c_- e^{-i\omega
    t}$; taking real and imaginary parts (with $c_\pm$ complex
    conjugates of each other, since $u$ is real-valued) recovers the
    familiar solution \eqref{eq:harmonic_oscillator_solution} written
    in terms of $\cos(\omega t)$ and $\sin(\omega t)$.
\end{example}


\section{Solving PDEs}

We now turn to PDEs of the general form \eqref{eq:second_order_pde}
from \Cref{ch:Lec1-intro}, and review a handful of classical techniques
for solving them in special cases; see \cite{strauss2007partial} for a
thorough treatment.

\subsection{Dealing with boundary conditions}

PDEs posed on a bounded domain $\Omega$ need boundary conditions on
$\partial\Omega$ to be well-posed, just like the BVPs of
\Cref{def:bvp}. The three classical types, using the normal derivative
notation $\partial u/\partial n$ from
\Cref{cor:greens-first-identity}, are:
\begin{itemize}
    \item \emph{Dirichlet}: $u = g$ on $\partial\Omega$, prescribing
    the value of the solution itself (e.g., a fixed temperature, as in
    \Cref{ex:heat_equation}).
    \item \emph{Neumann}: $\partial u/\partial n = g$ on
    $\partial\Omega$, prescribing the flux across the boundary (e.g.,
    $g=0$ models a perfectly insulated boundary with no heat flux).
    \item \emph{Robin (mixed)}: $au + b \, \partial u/\partial n = g$
    on $\partial\Omega$ for constants $a,b$, a weighted combination of
    the two (e.g., modeling heat loss proportional to the temperature
    difference with the surroundings).
\end{itemize}
Different boundary conditions imposed on the same PDE can lead to very
different solutions, so identifying the correct type from the physical
(or other) context is an important first step in solving any PDE.

\subsection{Reduction to ODEs}

The following four ansatzes all work by reducing a PDE to one or more
ODEs, which can then be attacked with the tools from the previous
section.

\paragraph{Separation of variables.}
We look for solutions of the special product form $u(t,x) =
T(t)X(x)$. Substituting into a linear PDE and dividing by $TX$
typically separates the equation into two sides depending on different
variables, which must therefore both equal the same separation
constant, giving one ODE per variable.

\begin{example}
    \label{ex:heat-separation}
    Consider the heat equation from \Cref{ex:heat_equation} with
    $u(t,x) = T(t)X(x)$. Substituting into $u_t = \alpha u_{xx}$ and
    dividing by $\alpha TX$ gives
    \begin{equation*}
        \frac{T'(t)}{\alpha T(t)} = \frac{X''(x)}{X(x)} = -\lambda,
    \end{equation*}
    for some constant $-\lambda$ (the sign is chosen for later
    convenience), since the left side depends only on $t$ and the right
    side only on $x$. This yields $T' = -\lambda \alpha T$ and $X'' =
    -\lambda X$, together with the boundary conditions $X(0)=X(L)=0$
    inherited from $u(t,0)=u(t,L)=0$. The latter is itself an
    eigenvalue problem for $X$, with solutions
    \begin{equation*}
        X_n(x) = \sin\left( \frac{n\pi x}{L} \right), \qquad \lambda_n =
        \left( \frac{n\pi}{L} \right)^2, \qquad n=1,2,\ldots,
    \end{equation*}
    and correspondingly $T_n(t) = e^{-\alpha \lambda_n t}$. Each
    $u_n(t,x) = X_n(x)T_n(t)$ is called a separated solution; we
    revisit this example below.
\end{example}

\paragraph{Similarity and self-similar solutions.}
Instead of a product ansatz, we look for a solution that depends on
the original variables only through a specific combination of them
(the similarity variable), reducing the number of independent
variables by one.

\begin{example}
    The heat equation $u_t = \alpha u_{xx}$ posed on all of $\R$
    (rather than a bounded interval) admits the self-similar solution
    \begin{equation*}
        u(t,x) = \frac{1}{\sqrt{4\pi\alpha t}} \exp\left( -\frac{x^2}{4
        \alpha t} \right), \qquad t > 0,
    \end{equation*}
    which depends on $(t,x)$ only through the similarity variable
    $\eta = x/\sqrt{4\alpha t}$, up to an overall prefactor; one can
    check directly by substitution that it solves the heat equation.
    This particular solution is called the \emph{heat kernel}, and we
    will see below that it plays the role of a Green's function for
    the heat equation.
\end{example}

\paragraph{Eigenfunction expansions.}
Separation of variables produces infinitely many separated solutions
$u_n$, one for each eigenvalue. Since the PDE is linear, any
(convergent) superposition $u = \sum_n c_n u_n$ is again a solution,
and by choosing the coefficients $c_n$ appropriately we can match a
prescribed initial (or boundary) condition that is not itself a single
eigenfunction.

\begin{example}
    Returning to \Cref{ex:heat-separation}, the general solution of the
    heat equation with zero Dirichlet boundary conditions is
    \begin{equation*}
        u(t,x) = \sum_{n=1}^\infty c_n \sin\left( \frac{n\pi x}{L}
        \right) e^{-\alpha (n\pi/L)^2 t},
    \end{equation*}
    where the $c_n$ are exactly the Fourier sine series coefficients of
    the initial condition $u_0(x) = u(0,x)$,
    \begin{equation*}
        c_n = \frac{2}{L} \int_0^L u_0(x) \sin\left( \frac{n\pi x}{L}
        \right) dx.
    \end{equation*}
    This is the analytic solution of \Cref{ex:heat_equation}, and it is
    instructive to compare it to the forward and backward Euler
    approximations we computed there numerically.
\end{example}

\paragraph{Traveling waves.}
We look for solutions of the form $u(t,x) = w(x-ct)$ for a wave speed
$c$ and profile function $w$, which propagate at constant speed $c$
without changing shape. Substituting this ansatz into the PDE and
using the chain rule reduces it to an ODE for $w$ in the single
variable $\xi = x-ct$.

\begin{example}
    The transport equation $u_t + cu_x = 0$ admits traveling wave
    solutions $u(t,x) = w(x-ct)$ for \emph{any} differentiable profile
    $w$: indeed $u_t = -cw'(x-ct)$ and $u_x = w'(x-ct)$, so $u_t + cu_x
    = 0$ identically. In particular, $u(t,x) = u_0(x-ct)$ solves the
    transport equation with initial condition $u(0,x) = u_0(x)$: the
    initial profile is simply transported rigidly to the right at
    speed $c$ (or to the left if $c<0$).
\end{example}

\subsection{Method of characteristics}

The method of characteristics solves first-order PDEs of the form
$u_t + c(t,x) u_x = f(t,x,u)$ by finding curves $(t,x(t))$ in the
$(t,x)$-plane, called characteristics, along which the PDE reduces to
an ODE in $t$ alone. Along a curve with $dx/dt = c(t,x)$, the chain
rule gives
\begin{equation*}
    \frac{d}{dt} u(t,x(t)) = u_t + \frac{dx}{dt} u_x = u_t + c(t,x) u_x
    = f(t,x(t),u(t,x(t))),
\end{equation*}
so $u$ along the characteristic solves the ODE $du/dt =
f(t,x(t),u)$. Solving the characteristic ODE $dx/dt = c(t,x)$ together
with this equation for $u$, and tracing back to $t=0$, gives the
solution everywhere the characteristics are well defined.

\begin{example}
    Consider the variable-speed transport equation $u_t + xu_x = 0$,
    $u(0,x) = u_0(x)$. The characteristics solve $dx/dt = x$, so $x(t)
    = x_0 e^t$ where $x_0 = x(0)$, and $u$ is constant along each
    characteristic since $f=0$. Solving for $x_0 = xe^{-t}$ and using
    $u(t,x(t)) = u(0,x_0) = u_0(x_0)$ gives
    \begin{equation*}
        u(t,x) = u_0(xe^{-t}).
    \end{equation*}
\end{example}

\subsection{Transformations}

The Fourier and Laplace transforms, as well as Green's functions,
extend naturally from the ODE setting of the previous section to PDEs,
by transforming only some of the independent variables.

\paragraph{Fourier.}
Taking the Fourier transform in $x$ of a PDE (leaving $t$
untransformed) turns spatial derivatives into multiplication by powers
of $i\omega$, reducing the PDE to a family of ODEs in $t$, one for each
frequency $\omega$. For instance, Fourier transforming the heat
equation $u_t = \alpha u_{xx}$ in $x$ gives $\hat u_t(t,\omega) =
-\alpha \omega^2 \hat u(t,\omega)$, an ODE in $t$ for each fixed
$\omega$ with solution $\hat u(t,\omega) = \hat u_0(\omega)
e^{-\alpha\omega^2 t}$; inverting the transform recovers exactly the
self-similar heat kernel solution seen above (convolved with the
initial data $u_0$).

\paragraph{Laplace.}
Similarly, taking the Laplace transform in $t$ eliminates time
derivatives (in favor of the initial data, via $\mathcal{L}\{u_t\} = sU
- u_0$) and reduces the PDE to a boundary value problem in the spatial
variable alone, for each value of the transform variable $s$; this is
a common strategy for PDEs on semi-infinite time intervals.

\paragraph{Green's functions and fundamental solutions.}
The Green's function idea from the previous section also extends to
PDEs: the \emph{fundamental solution} $\Phi$ of a linear PDE operator
$L$ solves $L\Phi = \delta$ in free space (i.e., with no boundary, or
with boundary conditions ``at infinity''), and the heat kernel from
above is exactly the fundamental solution of the heat operator $L =
\partial_t - \alpha \partial_x^2$. On a bounded domain with boundary
conditions, the Green's function of
\Cref{ex:greens-function-poisson} generalizes in the same way to PDE
boundary value problems such as Poisson's equation in higher
dimensions.

\bibliographystyle{plainnat}
\bibliography{references}